
générateurs de @@M0@@ au voisinage de @@M1@@ tel que le complexe de Koszul associé soit acyclique (il en existe si on suppose l'immersion régulière); soit @@M2@@ un autre système minimum de générateurs de @@M3@@; il existe un voisinage de @@M4@@ sur lequel @@M5@@ peut être engendré soit par les @@M6@@, soit par les @@M7@@, et par suite, il existe sur ce voisinage une matrice à coefficients dans @@M8@@ permettant d'exprimer les @@M9@@ en fonction des @@M10@@. Or les images des @@M11@@ (resp. @@M12@@) dans @@M13@@ forment une base au point @@M14@@, donc dans un voisinage de @@M15@@, et par suite cette matrice est inversible dans un voisinage de @@M16@@. Elle permet donc de définir un isomorphisme du complexe de Koszul défini par les @@M17@@ sur le complexe de Koszul défini par les @@M18@@, ce qui montre que le complexe de Koszul défini par les @@M19@@ est acyclique.

\textbf{1.4.3.} En supposant seulement @@M20@@ de type fini, et @@M21@@ localement libre de type fini, notons que tout système minimum de générateurs de @@M22@@ au voisinage d'un point @@M23@@ de @@M24@@ donne une base de @@M25@@ au voisinage de @@M26@@. Réciproquement, si on a au voisinage de @@M27@@ une base de @@M28@@, on peut la relever au voisinage de @@M29@@ en une famille de sections de @@M30@@, qui engendrent @@M31@@ par Nakayama, donc @@M32@@ dans un voisinage de @@M33@@ puisque @@M34@@ est de type fini; de plus ces générateurs sont en nombre minimum, puisqu'ils donnent une base de @@M35@@.

\textbf{Proposition 1.5.} (\uline{Comparer} EGA IV 19.1.5 (ii)). \uline{Soit} @@M36@@ \uline{une immersion}. \uline{Si} @@M37@@ \uline{est régulière}, \uline{il en est de même pour toute immersion déduite de} @@M38@@ \uline{par changement de base plat}. \uline{Pour que} @@M39@@ \uline{soit régulière}, \uline{il faut et il suffit qu'elle le soit localement pour la topologie fidèlement plate quasi-compacte}.

La stabilité par changement de base plat est évidente.

Réciproquement, soit @@M40@@ l'immersion déduite de @@M41@@ par un changement de
